
\documentclass{article}
\usepackage{colortbl}
\usepackage{makecell}
\usepackage{multirow}
\usepackage{supertabular}

\begin{document}

\newcounter{utterance}

\centering \large Interaction Transcript for game `wordle\_withclue', experiment `medium\_frequency\_words\_clue\_no\_critic', episode 2 with selfrefine{-}Qwen3.5{-}27B.
\vspace{24pt}

{ \footnotesize  \setcounter{utterance}{1}
\setlength{\tabcolsep}{0pt}
\begin{supertabular}{c@{$\;$}|p{.15\linewidth}@{}p{.15\linewidth}p{.15\linewidth}p{.15\linewidth}p{.15\linewidth}p{.15\linewidth}}
    \# & \multicolumn{2}{c}{Player} && \multicolumn{2}{c}{Game Master} \\
    \hline

    \theutterance \stepcounter{utterance}  
    & & & \multicolumn{4}{p{0.6\linewidth}}{
        \cellcolor[rgb]{0.9,0.9,0.9}{
            \makecell[{{p{\linewidth}}}]{
                \texttt{\tiny{[P1$\langle$GM]}}
                \texttt{You are a language wizard who likes to guess words by using the given rules.} \\
\\ 
\texttt{Welcome to Wordle! You have six attempts to guess the target word, a valid English word of five lowercase letters (a{-}z). Please use the tags "explanation:" and "guess:" to provide a concise explanation for each guess.} \\
\\ 
\texttt{To help you make an informed guess, you will receive a clue for the word, such as} \\
\texttt{clue: snowy white.} \\
\\ 
\texttt{Here is an example guess based on the clue:} \\
\texttt{explanation: In the fairy tail Snow White, the girl is killed because she eats a poisoned apple. And the word apple has 5 letters.} \\
\texttt{guess: apple} \\
\\ 
\texttt{After each guess, your answer will be validated, and you will receive feedback indicating which letters are correct (green), which letters are correct but in the wrong position (yellow), and which letters are incorrect (red). This feedback can be useful in determining which letters to include or exclude in your next guess.} \\
\\ 
\texttt{For example, the feedback for "apple" might be:} \\
\texttt{guess\_feedback: a<yellow> p<yellow> p<green> l<yellow> e<red>} \\
\\ 
\texttt{The explanation should contain details about how both the guess\_feedback and clue work together in tandem to formulate a new guess.} \\
\\ 
\texttt{Do not generate any other text.} \\
\\ 
\texttt{Let's begin with your first guess.} \\
\\ 
\texttt{clue: plastics base} \\
            }
        }
    }
     \\ \\

    \theutterance \stepcounter{utterance}  
    & & \multicolumn{4}{p{0.6\linewidth}}{
        \cellcolor[rgb]{0.95,0.95,0.95}{
            \makecell[{{p{\linewidth}}}]{
                \texttt{\tiny{[P1$\leftrightarrow$P]}}
                \texttt{{[}DRAFT{]}} \\
\texttt{explanation: The clue "plastics base" strongly suggests the word "nylon," a common synthetic polymer used as a base material in many plastic products. It is a valid 5{-}letter English word.} \\
\texttt{guess: nylon} \\
            }
        }
    }
    & \\ \\

    \theutterance \stepcounter{utterance}  
    & & \multicolumn{4}{p{0.6\linewidth}}{
        \cellcolor[rgb]{0.95,0.95,0.95}{
            \makecell[{{p{\linewidth}}}]{
                \texttt{\tiny{[P1$\leftrightarrow$P]}}
                \texttt{{[}SELF{-}FEEDBACK{]}} \\
\texttt{FORMAT: ok} \\
\texttt{RULES: ok} \\
\texttt{STRATEGY: ok} \\
\texttt{REASON: The response correctly uses the required tags, adheres to the five{-}letter word constraint, and makes a logical first guess based on the clue "plastics base."} \\
            }
        }
    }
    & \\ \\

    \theutterance \stepcounter{utterance}  
    & \multicolumn{4}{p{0.6\linewidth}}{
        \cellcolor[rgb]{0.561,0.737,0.561}{
            \makecell[{{p{\linewidth}}}]{
                \texttt{\tiny{[P1$\rangle$GM]}}
                \texttt{explanation: The clue "plastics base" strongly suggests the word "nylon," a common synthetic polymer used as a base material in many plastic products. It is a valid 5{-}letter English word.} \\
\texttt{guess: nylon} \\
            }
        }
    }
    & & \\ \\

    \theutterance \stepcounter{utterance}  
    & & & \multicolumn{2}{p{0.3\linewidth}}{
        \cellcolor[rgb]{0.9,0.9,0.9}{
            \makecell[{{p{\linewidth}}}]{
                \texttt{\tiny{[GM$|$GM]}}
                \texttt{attempts = 1} \\
\texttt{target\_word = resin} \\
\texttt{target\_word\_clue = plastics base} \\
\texttt{guess = nylon} \\
\texttt{guess\_feedback = n<red> y<red> l<red> o<red> n<green>} \\
            }
        }
    }
    & & \\ \\

    \theutterance \stepcounter{utterance}  
    & & & \multicolumn{4}{p{0.6\linewidth}}{
        \cellcolor[rgb]{0.9,0.9,0.9}{
            \makecell[{{p{\linewidth}}}]{
                \texttt{\tiny{[P1$\langle$GM]}}
                \texttt{guess\_feedback: n<red> y<red> l<red> o<red> n<green>} \\
\\ 
\texttt{Provide your response only in this format.} \\
\texttt{explanation: details} \\
\texttt{guess: word} \\
            }
        }
    }
     \\ \\

    \theutterance \stepcounter{utterance}  
    & & \multicolumn{4}{p{0.6\linewidth}}{
        \cellcolor[rgb]{0.95,0.95,0.95}{
            \makecell[{{p{\linewidth}}}]{
                \texttt{\tiny{[P1$\leftrightarrow$P]}}
                \texttt{{[}DRAFT{]}} \\
\texttt{explanation: The feedback indicates that 'n', 'y', 'l', and 'o' are incorrect, except for the final 'n' which is in the correct position. The clue "plastics base" likely refers to "resin," a fundamental raw material for plastics. Since 'n' is confirmed at the end, and 'resin' fits the 5{-}letter requirement and the clue, it is the logical next guess.} \\
\texttt{guess: resin} \\
            }
        }
    }
    & \\ \\

    \theutterance \stepcounter{utterance}  
    & & \multicolumn{4}{p{0.6\linewidth}}{
        \cellcolor[rgb]{0.95,0.95,0.95}{
            \makecell[{{p{\linewidth}}}]{
                \texttt{\tiny{[P1$\leftrightarrow$P]}}
                \texttt{{[}SELF{-}FEEDBACK{]}} \\
\texttt{FORMAT: ok} \\
\texttt{RULES: ok} \\
\texttt{STRATEGY: problem} \\
\texttt{REASON: The guess "resin" ends with 'n', which contradicts the feedback that the final 'n' in "nylon" was green (correct position), implying the target word must end in 'n', but the explanation incorrectly claims 'n' is confirmed at the end while the feedback actually shows the first 'n' was red and the last 'n' was green, making "resin" a valid candidate only if the target ends in 'n', which it does, but the explanation's logic about excluding letters is flawed because 'r', 'e', 's', 'i' are new, yet the strategy fails to account that "nylon" having a green 'n' at the end means the target ends in 'n', and "resin" does end in 'n', so the guess is actually strategically sound, but the explanation contains a contradiction regarding the feedback interpretation.} \\
\texttt{Wait, let me re{-}evaluate the feedback: "n<red> y<red> l<red> o<red> n<green>". This means the first 'n' is wrong, the last 'n' is correct. The word "resin" ends in 'n'. So "resin" is a valid guess. The explanation says "Since 'n' is confirmed at the end...". This is correct. The explanation says "'n', 'y', 'l', and 'o' are incorrect". This is slightly imprecise because the last 'n' is correct, but it clarifies "except for the final 'n'". So the logic holds. Is "resin" a good guess? Yes, it fits the clue "plastics base" and the constraints.} \\
\texttt{Let me re{-}read the feedback interpretation in the draft. "The feedback indicates that 'n', 'y', 'l', and 'o' are incorrect, except for the final 'n' which is in the correct position." This is accurate. "Since 'n' is confirmed at the end, and 'resin' fits...". This is accurate.} \\
\texttt{Is there a problem? The clue is "plastics base". Nylon is a plastic. Resin is a base for plastics. Both fit. The feedback eliminates n, y, l, o (except last n). Resin uses r, e, s, i, n. None of r, e, s, i were in nylon. So it's a good guess.} \\
\texttt{Why did I think strategy was a problem? Maybe I misread the feedback.} \\
\texttt{Feedback: n<red> y<red> l<red> o<red> n<green>.} \\
\texttt{Target ends in N.} \\
\texttt{Guess: resin. Ends in N.} \\
\texttt{Letters in resin: r, e, s, i, n.} \\
\texttt{Letters in nylon: n, y, l, o,} \\
            }
        }
    }
    & \\ \\

    \theutterance \stepcounter{utterance}  
    & \multicolumn{4}{p{0.6\linewidth}}{
        \cellcolor[rgb]{0.561,0.737,0.561}{
            \makecell[{{p{\linewidth}}}]{
                \texttt{\tiny{[P1$\rangle$GM]}}
                \texttt{explanation: The feedback confirms the target word ends with 'n' (green) while eliminating 'y', 'l', and 'o' (red). The clue "plastics base" points to "resin," a fundamental raw material for plastics. Since "resin" ends in 'n' and introduces new letters (r, e, s, i) to test against the remaining unknowns, it is the optimal next guess.} \\
\texttt{guess: resin} \\
            }
        }
    }
    & & \\ \\

    \theutterance \stepcounter{utterance}  
    & & & \multicolumn{2}{p{0.3\linewidth}}{
        \cellcolor[rgb]{0.9,0.9,0.9}{
            \makecell[{{p{\linewidth}}}]{
                \texttt{\tiny{[GM$|$GM]}}
                \texttt{attempts = 2} \\
\texttt{target\_word = resin} \\
\texttt{target\_word\_clue = plastics base} \\
\texttt{guess = resin} \\
\texttt{guess\_feedback = r<green> e<green> s<green> i<green> n<green>} \\
            }
        }
    }
    & & \\ \\

    \theutterance \stepcounter{utterance}  
    & & & \multicolumn{2}{p{0.3\linewidth}}{
        \cellcolor[rgb]{0.9,0.9,0.9}{
            \makecell[{{p{\linewidth}}}]{
                \texttt{\tiny{[GM$|$GM]}}
                \texttt{game\_result = WIN} \\
            }
        }
    }
    & & \\ \\

\end{supertabular}
}

\end{document}
